\section{Semirreta}

A semirreta \texttt{(a, b)} começa em $a$ e segue para sempre na direção de $b$. Se $a = b$, ela é só o ponto $a$.
\texttt{on\_ray(a, b, p)} diz se $p$ está nela: na reta (\texttt{orient} $= 0$) e não atrás de $a$
(\texttt{dot(b - a, p - a)} $\ge 0$).

\subsection*{Distância de um ponto}

\begin{drawing}
  \fill[fig_fill] (0,-1.6) rectangle (6.5,2.2);
  \draw[gray, dashed] (0,-1.6) -- (0,2.2);
  \draw[->, thick] (0,0) -- (6.5,0);
  \fill (0,0) circle (1.5pt) node[below left] {$a$};
  \fill (2,0) circle (1.5pt) node[below] {$b$};
  \fill (4,1.6) circle (1.5pt) node[above] {$p$};
  \draw[fig_highlight, thick] (4,1.6) -- (4,0);
  \draw (4,0.25) -- (4.25,0.25) -- (4.25,0);
  \fill (-2,1.3) circle (1.5pt) node[above] {$p'$};
  \draw[fig_highlight, thick] (-2,1.3) -- (0,0);
  \node[fig_accent] at (5.2,-1.1) {dot $> 0$};
\end{drawing}

\texttt{ray\_point(a, b, p)}: se $p$ está à frente de $a$ (região azul, \texttt{dot(b - a, p - a)} $> 0$), a
distância é até a reta, pela perpendicular. Se está atrás, como $p'$, o ponto mais perto da semirreta é a própria
origem $a$.

\subsection*{Distância entre duas semirretas}

\begin{drawing}
  \draw[->, thick] (0,0) -- (3,2);
  \draw[->, thick] (0,2.2) -- (3,0);
  \fill (0,0) circle (1.5pt) node[left] {$a$};
  \fill (0,2.2) circle (1.5pt) node[left] {$c$};
  \fill[fig_highlight] (1.571,1.048) circle (2pt);
  \node at (1.5,-0.6) {se cruzam: 0};
  \draw[->, thick] (4.5,0) -- (8,0);
  \fill (4.5,0) circle (1.5pt) node[below] {$a$};
  \draw[->, thick] (6,1.8) -- (7.8,2.4);
  \fill (6,1.8) circle (1.5pt) node[left] {$c$};
  \draw[fig_highlight, thick] (6,1.8) -- (6,0);
  \node at (6.2,-0.6) {não se cruzam};
\end{drawing}

\texttt{ray\_ray(a, b, c, d)}: as semirretas se encontram em $a + t \cdot v = c + s \cdot w$, com $v = b - a$ e
$w = d - c$. Se não são paralelas (\texttt{cross(v, w)} $\ne 0$), elas se cruzam quando $t \ge 0$ e $s \ge 0$. Como
\[ t = \frac{\text{cross}(c - a, w)}{\text{cross}(v, w)}, \qquad s = \frac{\text{cross}(c - a, v)}{\text{cross}(v, w)}, \]
basta comparar o sinal dos numeradores com o do denominador, sem dividir: o teste é exato com \texttt{T} inteiro.
Se não se cruzam (ou são paralelas), os pontos mais próximos incluem uma das origens, então a resposta é o menor entre
\texttt{ray\_point(a, b, c)} e \texttt{ray\_point(c, d, a)}.
